% Separate post-CP25 proof. Four complete actual-file audits PASS. Ordinary hereditary count on arbitrary punctured covers.
\subsection{Degree-linear ordinary representation growth on arbitrary punctured covers}

\begin{theorem}[Proved: hereditary ordinary representation growth]
\label{D:arbitrary-cover-ordinary-hereditary-growth}
Let $(R,\mathfrak m,k)$ be a complete excellent normal local
surface with algebraically closed coefficient and residue field
$k$ of characteristic $p>0$. Put
$U_R=\operatorname{Spec}R\setminus\{\mathfrak m\}$.
There is a constant $A_R$ with the following property.
For every connected finite etale cover $U_A\to U_R$ of
degree $t$, every $q=p^f$, and every $r\ge1$,
\[
 \left|\operatorname{Hom}_{\rm cont}
       (\pi_1(U_A),\operatorname{GL}_r(\mathbb F_q))
          /\operatorname{GL}_r(\mathbb F_q)\right|
                         \le q^{A_Rtr^2}.                \tag{D.HU.1}
\]
There is no tameness condition on the finite cover.
The constant is independent of its degree, its ramification
along a resolution, $r$, and $f$.

More precisely, after choosing one of at most
\[
              q^{c_0t}\exp(c_1ft\log(r+1))               \tag{D.HU.2}
\]
fixed ordered determinant tuples, all the relevant ordinary
unit objects are covered by a single open subset of an affine
polynomial zero scheme with at most $c_2ftr^2$ coordinates.
Its equations have degree at most a constant $D_R$ independent
of $t,r,f$. Every point of this open parameter scheme produces
an actual formal, and on finite-type coefficient bases actual
algebraic, Frobenius unit object on $U_A$.
Its punctured unit classes have open and closed images.
The parameter families need not be relatively canonical,
coflat, or flat families of maximal Cohen--Macaulay modules.
\end{theorem}

\begin{proof}
The finite normalization $A$ is a complete normal local domain
with residue field $k$: it is finite over the complete henselian
base, and connectedness of the punctured cover selects one
normal local factor. Use the one fixed henselian algebraization
and all-cover geometric budget results to choose the specified
regular projective model $\pi:Y\to\operatorname{Spec}A$.
On the fixed base model choose the separating parameter pencil
inside the cover-uniform Noether-different ideal. Apply
Proposition~\ref{D:uniform-comparable-reduction-cycle}.
It provides actual cycles
\[
 C=f_A^*C_0,\qquad Z=\alpha C-\beta R_{f_A},\qquad
                C\le Z\le\alpha C,                     \tag{D.HU.3}
\]
where $\alpha,\beta$ and the base data depend only on $R$.
The number of exceptional components and $p_g(A)$ are $O_R(t)$
on this same model. In addition,
\[
 \begin{gathered}
 (x^j,y^j)\mathcal O_Y=\mathcal O_Y(-jC),\qquad
 -C^2,\ -Z\cdot C,\ -Z^2,\ |K_Y\cdot Z|=O_R(t),\\
 H^1(Z,\mathcal O_Z(-jZ))=0\quad(j\ge1),\qquad
 \mathcal O_Y(-C),\ \mathcal O_Y(-Z)
          \text{ are globally generated}.
 \end{gathered}                                        \tag{D.HU.4}
\]
Here generation and the conormal vanishings follow from the
all-subcycle inequalities in that proposition and the checked
uniform-generation theorem. They apply to every fixed multiple
of $Z$. Generation of $\mathcal O_Y(-C)$ also follows directly
from the parameter pencil.

For clarity, the scalar budgets on any fixed $\ell C$ are
simultaneous with (D.HU.4). Indeed
$K_Y=f_A^*K_{Y_0}+R_{f_A}$, with
$0\le R_{f_A}\le\lambda C$ and $C$ anti-nef. Projection
therefore gives $|K_Y\cdot C|=O_R(t)$, as well as
$-C^2=O_R(t)$. Coherent $H^2$ vanishes over the affine
surface, so $H^1(Y,\mathcal O_Y)$ surjects onto
$H^1(\ell C,\mathcal O_{\ell C})$. Adjunction now gives
\[
 h^0(\mathcal O_{\ell C})+h^1(\mathcal O_{\ell C})
                              =O_{R,\ell}(t).            \tag{D.HU.5}
\]
No additional blowup of this chosen model is made.

\emph{Geometric determinant lists and canonical input objects.}
A continuous $\mathbb F_q$ representation is equivalently an
ordinary $q$-unit bundle on $U_A$, or a length-$f$ cycle of
ordinary $p$-unit bundles. At phase $a$ take its reflexive
extension $M_a$ and its canonical full pullback
$F_a=(\pi^*M_a)^{\vee\vee}$. The actual canonical arrows
satisfy
\[
 \phi_a:F_Y^*F_a\longrightarrow F_{a+1},\qquad
 \phi_a^{-1}:F_{a+1}\longrightarrow(F_Y^*F_a)(pC).
                                                               \tag{D.HU.6}
\]
These are regular maps to the indicated bundles. They are
the canonical Frobenius sandwich supplied by the first-power
Noether factorization and the parameter pencil, as in
Proposition~\ref{D:canonical-family-unit-Isom-openness}.
For these geometric input objects no relative fullness
assertion is needed.

The proof of Lemma~\ref{D:hereditary-ordinary-unit-determinants}
applies on this model. We recall why its constants remain
unchanged when the generating cycle is now $Z$.
There are at most $q^{c_0t}$ initial punctured determinant
characters, because its bound
$|H^1(U_A,\boldsymbol\mu_{q-1})|
 \le(q-1)^{2p_g(A)+v}$ uses only the geometric genus and the
number $v=O_R(t)$ of exceptional components. All later
punctured determinants are the prescribed $p$-power phases
of that one character. Generic evaluation of a full bundle
gives nonnegative integral exceptional determinant degrees.
Pairing the cyclic determinant recurrence of (D.HU.6) with
the anti-nef cycle $C$ gives
\[
 0\le c_1(F_a)\cdot C
              \le\frac{rp}{p-1}(-C^2)=O_R(tr).          \tag{D.HU.7}
\]
Every coefficient of $C$ is a positive integer. Thus the
unweighted sum of these degrees is also $O_R(tr)$.
For a fixed punctured determinant, the negative definite
intersection matrix makes its exceptional multidegree
determine its extension line on $Y$ uniquely. The same
binomial enumeration as in the cited lemma gives at most
$\exp(c_1t\log(r+1))$ determinant lines per phase.
This proves (D.HU.2) for the ordered global lines
\[
                    L_a=\det(F_a(-Z)).                 \tag{D.HU.8}
\]
These are choices of actual lines on $Y$, not moving Picard
parameters. The cycle $Z$ is fixed before the representation.
By (D.HU.3) and the comparable-cycle generation theorem,
every $F_a(-Z)$, and hence every $L_a$, is globally generated.

\emph{One fixed finite input level.}
Choose once
\[
       b=p,\qquad N\ge\max\{3,\alpha\},\qquad L=N+p,
       \qquad D=LC.                                    \tag{D.HU.9}
\]
Then $pN-p>N$ and $Z\subset NC\subset D$ as actual
Cartier subschemes. Fix one ordered line tuple (D.HU.8).
On $D$ the nonnegative degrees of its restrictions satisfy
\[
 \deg_D(L_a)=L\bigl(c_1(F_a)\cdot C-rZ\cdot C\bigr)
                                                =O_R(tr).
                                                               \tag{D.HU.10}
\]
In conjunction with (D.HU.5), the regular-section exact
sequence for a generated line gives
$h^0(D,L_a)+h^1(D,L_a)=O_R(tr)$.

Use the actual quotient presentation of
Proposition~\ref{D:hereditary-ordinary-unit-finite-cycle-Hom}:
on the generating-tuple open locus in
$H^0(D,L_a)^{r+1}$ there is a locally split exact sequence
\[
 0\longrightarrow L_a^{-1}\xrightarrow{s_a}
       \mathcal O_D^{r+1}\longrightarrow V_a
       \longrightarrow0,
       \qquad E_a=V_a(Z).                               \tag{D.HU.11}
\]
The universal quotients are actual vector bundles after
every coefficient change, and their determinants are the
specified $L_a$. Every geometric canonical bundle
$F_a(-Z)|_D$ occurs, because a generated rank-$r$ bundle
on this proper Cohen--Macaulay curve has $r+1$ general
generators. The total number of object coordinates is
$O_R(ftr^2)$.

\emph{Exact arrow coordinates and their degree bound.}
We must encode the finite data of
Proposition~\ref{D:finite-Frobenius-object-effectivity},
not merely independent extensions of objects. Introduce
actual maps on $D$
\[
 \alpha_a:F^*E_a\to E_{a+1},\qquad
 \beta_a:E_{a+1}\to(F^*E_a)(pC),                         \tag{D.HU.12}
\]
and impose both exact equalities
\[
 \alpha_a(pC)\beta_a=\iota_{E_{a+1}},\qquad
 \beta_a\alpha_a=\iota_{F^*E_a}.                         \tag{D.HU.13}
\]
These are the two Cartier inclusion maps on $D$; no
inversion of the Cartier equation on $D$ occurs.

Here is a full check of the coordinate bound for these
particular arrows. Frobenius pullback of (D.HU.11) is
exact over every ring because the sequence is locally
split. Its line term is $L_a^{-p}$, and its differential
is the $p$-power of the tuple $s_a$. The two residual
line twists in (D.HU.12) are respectively
\[
                    (1-p)Z,\qquad(p-1)Z+pC.              \tag{D.HU.14}
\]
The source, target and product Hom complexes consequently
involve only bounded powers and products of
$L_a$, $L_{a+1}$, $\mathcal O_D(-Z)$,
$\mathcal O_D(-C)$ and their inverses. These four kinds
of positive lines are generated, with degrees $O_R(tr)$,
$O_R(tr)$, $O_R(t)$ and $O_R(t)$ respectively.
Every needed line can therefore be written $P\otimes Q^{-1}$,
where $P,Q$ are generated and
$\deg P+\deg Q=O_R(tr)$.
Regular sections on the full nonreduced Cohen--Macaulay
curve give
\[
 h^0(D,PQ^{-1})+h^1(D,PQ^{-1})
 \le h^0(\mathcal O_D)+h^1(\mathcal O_D)+\deg P+\deg Q
                                  =O_R(tr).             \tag{D.HU.15}
\]

Apply the fixed two-affine Cech construction of
Proposition~\ref{D:fixed-determinant-finite-Hom} to these
fixed line terms. The Hom between the two-term quotient
complexes has four types of terms: a fixed-cycle line term with
$(r+1)^2$ copies and cohomology $O_R(t)$, two line terms with $r+1$ copies, and
one line term with one copy. By (D.HU.5) and (D.HU.15),
their total finite cohomology dimension is $O_R(tr^2)$.
All Cech splittings are chosen once over $k$, with no
variable cohomology bases. The transferred differential
and the inclusion/projection maps are the exact
all-ring polynomial maps of that proposition.
Substitution of the Frobenius tuples raises object
coefficients only to their $p$-th powers; the internal
complex width is still two. Thus the degrees of these
maps and of composition are bounded by a constant
depending on $p$ alone.

A degree-zero cycle in each finite Hom complex represents
an actual map (D.HU.12). Impose its cycle equation.
For (D.HU.13), use the exact transferred composition and
one degree-minus-one boundary witness for each equality.
This is equivalent over every coefficient ring to equality
of the actual maps, even when cohomology is not flat.
There are a fixed number of such complexes, maps and
length-two relations per phase. The combined coordinates
and scalar equations are therefore $O_R(ftr^2)$ and
have a uniform finite degree bound $D_R$.
The chosen lines, Cech splittings, Cartier sections and
all multiplication coefficients are fixed background
constants; they may depend on the cover and line tuple,
but they introduce no variable coordinates or degree cost.

Let $X_{\boldsymbol L}$ be the affine zero scheme of
these polynomial equations in the stated coordinates.
Let $S_{\boldsymbol L}$ be its open generating-tuple
locus. This is one open subset of one polynomial zero
scheme, with no uncontrolled finite union of charts.
The properness defining the open quotient locus is used
only to obtain an open subset; its equations do not enter
the component estimate.

\emph{Actual families and the class comparison.}
Every point of $S_{\boldsymbol L}$ is actual finite
bounded-pole Frobenius data (D.HU.11)--(D.HU.13).
Proposition~\ref{D:finite-Frobenius-object-effectivity}
produces its formal unit family and retains the specified
marking on $NC$. On affine parameter opens the family
algebraizes over the proper Noetherian completed base.
On overlaps the unique integral unit comparison inducing
the specified identity $N$-jet satisfies its cocycle
identity. Thus these local constructions glue, and no
parameter-cover cardinality is charged in the estimate.

Because $Z\subset NC$, the constructed family retains
the actual generating quotient
$\mathcal O_{Z}^{r+1}\to F_a(-Z)|_Z$ from (D.HU.11),
and its specified determinant line $L_a|_Z$.
Lemma~\ref{D:weak-cycle-generated-determinant-lifting}
therefore lifts these generators formally and makes
$\det(F_a(-Z))$ the fixed global line $L_a$, locally on
the parameter scheme. The resulting objects need not
be canonical full pullbacks. That property is unnecessary:
Proposition~\ref{D:generated-fixed-determinant-unit-openness}
applies to the actual family and proves that its punctured
unit classes have disjoint open and closed images.
The comparison pole bounds may depend on the fixed
reference class and rank. Only qualitative finite
presentation is used there; no parameter of its Isom
scheme is added to $X_{\boldsymbol L}$.
Consequently
\[
 \#\{\text{punctured unit classes represented by }
                       S_{\boldsymbol L}\}
          \le\#\operatorname{Irr}(S_{\boldsymbol L})
          \le\#\operatorname{Irr}(X_{\boldsymbol L}).    \tag{D.HU.16}
\]

Every original representation occurs in these families.
Indeed its canonical objects and actual arrows restrict
to (D.HU.11)--(D.HU.13) for one chosen determinant tuple.
The true formal object is one permissible choice of local
matrix lifts in the object-effectivity construction.
Its choice-independence supplies an actual integral unit
isomorphism to the constructed object, retaining the
prescribed $N$-jet. Thus coverage is coverage of the
original punctured unit class, not just of its finite
truncation.

\emph{Counting components.}
The affine component bound of
Proposition~\ref{D:weighted-affine-component-bound}, with
all variable weights equal to one, gives, with
$M=\dim(\text{ambient affine space})\le c_2ftr^2$,
\[
 \#\operatorname{Irr}(X_{\boldsymbol L})
            \le(M+1)\max\{2,D_R\}^{M}
            \le p^{c_3ftr^2}.                          \tag{D.HU.17}
\]
It applies to arbitrary rectangular systems and counts
positive-dimensional components, so no gauge quotient or
zero-dimensional assumption has been inserted.
Multiplying this bound by (D.HU.2), and using
$\log(r+1)=O(r^2)$ and $\log q=f\log p$, gives
\[
 q^{c_0t}\exp(c_1ft\log(r+1))\,p^{c_3ftr^2}
                                      \le q^{A_Rtr^2}.
\]
Ordinary $q$-unit classes on $U_A$ are exactly the
conjugacy classes of continuous $\mathbb F_q$
representations. This proves (D.HU.1).
\end{proof}

\begin{remark}[Scope of the hereditary conclusion]
The theorem counts ordinary defining-characteristic
linear representations on arbitrary finite punctured
covers. No projective determinant enumeration, twisted
gerbe atlas, factorial bound for unrestricted alternating
quotients, or full profinite finite-generation conclusion
is inferred. The proof avoids coflat rank strata by
constructing all bounded-pole unit objects represented
by its finite equations and using generated
fixed-determinant comparison families.
\end{remark}
